\section{Summary}
\label{sec:design-summary}

This chapter developed and justified the architecture of the Reading the Reader platform, the central contribution of the thesis.
It first isolated the architecturally significant requirements as seven design drivers and adopted an inward-dependency, ports-and-adapters style as the structural response to the dominant driver of modular separability.
On that foundation it decomposed the platform's adaptive behaviour into four independently replaceable concerns (sensing, eye-movement analysis, decision, and intervention), each hidden behind a uniform contract and resolved through a registry, which is the modularity that RQ1 requires.
It set these concerns in motion as a real-time adaptive loop that reconciles modular separation with the responsiveness and robustness drivers, by confining the modular indirection to a lower-frequency control path while keeping the per-sample sensing path thin, and that commits interventions at controlled boundaries so that adaptation preserves the reader's position, as RQ2 and RQ3 require.
It described the single authoritative, schema-versioned session record that makes a session reproducible from the record alone and demonstrable through replay, the experimentability and auditable record that RQ4 examines.
Finally, it showed how the architecture is opened for extension, in-process through additive registration and out-of-process through a uniform provider seam, and gathered the design decisions and their alternatives into a single register.
The architecture established here is the specification that \cref{ch:implementation} realises in code, where selected implementation highlights, the build, and the development workflow are described.
